\documentclass[12pt,a4paper]{article}

\usepackage[T1]{fontenc}
\usepackage[utf8]{inputenc}
\usepackage[brazil]{babel}
\usepackage{lmodern}
\usepackage{geometry}
\usepackage{hyperref}
\usepackage{enumitem}
\usepackage{array}

\geometry{margin=2.5cm}
\hypersetup{
  colorlinks=true,
  linkcolor=blue,
  urlcolor=blue
}

\title{Plano de Aula: Aula 8 -- Planejamento e Gerenciamento de Software}
\date{}

\begin{document}
\maketitle

\noindent
\textbf{Disciplina:} Engenharia de Software (Curso Técnico em Desenvolvimento de Sistemas)\\
\textbf{Duração:} 2 horas (120 minutos)\\
\textbf{Modalidade:} Presencial\\
\textbf{Materiais da aula:} \href{aula-08.html}{slides (HTML)} \quad \href{aula-08.pdf}{ementa (PDF)}\\
\textbf{Livro-base (referência):} \href{../99-Materiais/Engenharia%20de%20software%20projetos%20e%20processos%20by%20Wilson%20de%20P%C3%A1dua%20Paula%20Filho%20.epub.pdf}{Engenharia de software: projetos e processos (PDF)}

\section{Competências e Habilidades}
\textbf{Competências}
\begin{itemize}[leftmargin=*,itemsep=0.25em]
  \item Aprender abordagens para planejamento e gerenciamento de projetos de software.
  \item Utilizar metodologias para organizar entregas e lidar com riscos.
\end{itemize}

\textbf{Habilidades}
\begin{itemize}[leftmargin=*,itemsep=0.25em]
  \item Planejar entregas com base em escopo, prioridades e dependências.
  \item Estimar esforço considerando incerteza (faixas).
  \item Identificar, priorizar e mitigar riscos.
  \item Controlar mudanças e acompanhar progresso.
\end{itemize}

\section{Objetivos da Aula}
\begin{itemize}[leftmargin=*,itemsep=0.25em]
  \item Entender por que projetos falham e como gerenciamento reduz riscos.
  \item Construir um plano mínimo de projeto (escopo, entregas, papéis, riscos, critérios).
  \item Estimar tarefas com faixas (otimista/provável/pessimista).
  \item Definir e aplicar uma matriz simples de risco.
  \item Compreender controle de mudanças e impacto.
\end{itemize}

\section{Assuntos Abordados no Dia}
\begin{itemize}[leftmargin=*,itemsep=0.25em]
  \item Plano mínimo de projeto: artefatos e decisões.
  \item Estimativas e incerteza: por que não existe precisão perfeita.
  \item Gestão de risco: exemplos, mitigação e matriz impacto x probabilidade.
  \item Controle de mudanças: registrar, analisar impacto, priorizar e comunicar.
  \item Atividade: plano mínimo para o ``app organizador''.
\end{itemize}

\section{Metodologia}
Exposição dialogada com exemplos do cotidiano e do ``app organizador''. Atividade prática em grupo para construir um plano mínimo (backlog priorizado, tarefas, estimativas, riscos, critérios de pronto) e discutir impacto de mudanças. O professor conduz uma revisão guiada e destaca sinais de risco.

\section{Materiais e Recursos}
\begin{itemize}[leftmargin=*,itemsep=0.25em]
  \item Quadro e marcadores.
  \item Folhas A4 para plano e matriz de risco (ou planilha, se disponível).
  \item Post-its (opcional) para backlog e tarefas.
  \item Slides: \href{aula-08.html}{aula-08.html} e ementa: \href{aula-08.pdf}{aula-08.pdf}.
  \item Livro-base (PDF): \href{../99-Materiais/Engenharia%20de%20software%20projetos%20e%20processos%20by%20Wilson%20de%20P%C3%A1dua%20Paula%20Filho%20.epub.pdf}{Engenharia de software: projetos e processos}.
  \item Vídeos complementares (links nos slides).
\end{itemize}

\section{Cronograma e Descrição Detalhada (120 Minutos)}
\subsection*{Parte 1: Abertura e motivação (10 min)}
\begin{itemize}[leftmargin=*,itemsep=0.35em]
  \item Retomar testes (Aula 7): qualidade depende de processo e gestão.
  \item Apresentar objetivo: reduzir surpresa no final com planejamento e risco.
\end{itemize}

\subsection*{Parte 2: Plano mínimo e visibilidade (25 min)}
\begin{itemize}[leftmargin=*,itemsep=0.35em]
  \item Apresentar componentes do plano: escopo, entregas, papéis, riscos, critérios.
  \item Discussão: o que acontece quando o trabalho não é visível.
\end{itemize}

\subsection*{Parte 3: Estimativa com faixas (20 min)}
\begin{itemize}[leftmargin=*,itemsep=0.35em]
  \item Separar tamanho de esforço e trabalhar com incerteza.
  \item Exercício rápido: estimar uma tarefa em três pontos (O/P/P).
\end{itemize}

\subsection*{Parte 4: Gestão de risco e matriz (25 min)}
\begin{itemize}[leftmargin=*,itemsep=0.35em]
  \item Definir risco e exemplos: requisitos, integração, tecnologia, pessoas.
  \item Montar matriz impacto x probabilidade e discutir prioridades de ação.
  \item Estratégias: mitigar, transferir, aceitar, evitar.
\end{itemize}

\subsection*{Parte 5: Controle de mudanças (15 min)}
\begin{itemize}[leftmargin=*,itemsep=0.35em]
  \item Fluxo simples: registrar, impacto, priorizar, atualizar plano e comunicar.
  \item Exemplos: mudança de requisito e impacto em prazo/qualidade.
\end{itemize}

\subsection*{Parte 6: Atividade prática (20 min)}
\begin{itemize}[leftmargin=*,itemsep=0.35em]
  \item Em grupos:
  \begin{itemize}[leftmargin=*,itemsep=0.25em]
    \item Escolher 6 itens de backlog e priorizar.
    \item Quebrar em tarefas e estimar com faixas.
    \item Listar 5 riscos e 1 mitigação por risco.
    \item Definir 3 critérios de ``pronto''.
  \end{itemize}
  \item Socialização rápida (2 grupos) e feedback do professor.
\end{itemize}

\subsection*{Parte 7: Fechamento (5 min)}
\begin{itemize}[leftmargin=*,itemsep=0.35em]
  \item Ponte para Aula 9: qualidade de produto e processo, métricas e melhoria contínua.
\end{itemize}

\section{Avaliação}
\begin{itemize}[leftmargin=*,itemsep=0.25em]
  \item Entrega do plano mínimo (com riscos e critérios).
  \item Coerência e justificativas (priorização e mitigação).
\end{itemize}

\section{Referências}
\begin{itemize}[leftmargin=*,itemsep=0.25em]
  \item \href{../99-Materiais/Engenharia%20de%20software%20projetos%20e%20processos%20by%20Wilson%20de%20P%C3%A1dua%20Paula%20Filho%20.epub.pdf}{PAULA FILHO, Wilson de Pádua. Engenharia de software: projetos e processos. 4. ed. 2019.}
\end{itemize}

\end{document}
